\setcounter{subsection}{0}
\section{DESCRIPCIÓN GENERAL}
%Aunque similar a la propuesta, no debería ser igual, se supone que debe actualizarse con nuevas referencias. Además, debe estar escrito en pasado, describiendo cuál era el problema y cómo se resolvió. 

\subsection{Oportunidad, Problema}
%Máximo 3 páginas.
%Presenta las características, contexto, situación o aspectos de interés de la oportunidad o pro-blema seleccionado. Los estudiantes pueden partir de un problema global y multidisciplinar (por ejemplo, uno de los temas planteados en la misión de la Pontificia Universidad Javeriana) y continuar con el problema resuelto o la oportunidad que se ha aprovechado con el proyecto. 
%Los puntos que podrían desarrollarse en este párrafo son similares a los de la propuesta. (Re-cuerda escribir en pasado, ya que se refiere al proyecto ya terminado): 


\subsubsection{Contexto del Problema}
%Describe las características, el contexto, la situación o el aspecto interesante del problema. Se recomienda explicar desde los conceptos más generales hasta los más específicos. La escritura debe ser coherente. 

\subsubsection{Formulación del Problema}
%Basándose en la información de la sección anterior, esta sección detalla el problema específico de esta propuesta y justifica su relevancia.

\subsubsection{Propuesta de Solución}
%Esta sección debe incluir una descripción general de la solución propuesta al problema anterior. Esta sección también debe indicar el área de la propuesta (por ejemplo, desarrollo de software, ingeniería de sistemas, etc.)

\subsubsection{Justificación de la Solución}
%En esta sección, los estudiantes explican por qué la solución es adecuada para abordar el pro-blema

\subsection{Descripción del Proyecto}
%Esta sección presenta la descripción general del proyecto. Cubre los aspectos más importantes para que el lector entienda qué se pretendía hacer.

\subsubsection{Objetivo General}
%Escribe aquí el objetivo general aprobado para el proyecto.

\subsubsection{Objetivos Específicos}
%Escribe aqui los objetivos específicos que, juntos, ayuden a cumplir el objetivo general.

\subsubsection{Marco metodológico, entregables, Estándares y Justificación}
%Describe la metodología, así como documentos y otros productos que se entregan en este pro-yecto final de grado. La información debe resumirse en la siguiente tabla:


\begin{table}[ht]
\renewcommand{\arraystretch}{1}
\centering
\begin{tabular}{|c|c|c|}
\hline
\textbf{Entregable} &
  \textbf{Estándares} &
  \textbf{Justificación} \\ \hline

\begin{tabular}[c]{@{}c@{}}   \\   \end{tabular} &   
   &
  \begin{tabular}[c]{@{}c@{}}  \end{tabular} \\ \hline
\begin{tabular}[c]{@{}c@{}}  \end{tabular} &  
   &
  \begin{tabular}[c]{@{}c@{}}  \end{tabular} \\ \hline
\end{tabular}
\caption{Entregables y Estándares Asociados}
\label{tab:entregables}
\end{table}